%%%%%%%%%%%%%%%%%%%%%%%%%%%%%%%%%%%%%%%%%%%%%%%%%%%%%%%%%%%%
% 10.14 DECISION THEORY
% The Composition of Advanced Mathematics
%%%%%%%%%%%%%%%%%%%%%%%%%%%%%%%%%%%%%%%%%%%%%%%%%%%%%%%%%%%%

\section{Decision Theory}
\label{sec:decision-theory}

\prerequisite{
Probability, random variables, expectation, optimization fundamentals,
statistics, Bayesian methods, stochastic optimization, and
multiobjective optimization.
}

\learningobjectives{
By the end of this section, the reader should be able to:
\begin{itemize}
    \item define a decision problem;
    \item distinguish decisions, states of nature, outcomes, and payoffs;
    \item construct payoff and loss tables;
    \item distinguish decisions under certainty, risk, and uncertainty;
    \item use expected monetary value;
    \item use expected loss;
    \item construct regret tables;
    \item apply minimax regret;
    \item apply maximax and maximin criteria;
    \item apply the Hurwicz criterion;
    \item apply the Laplace principle;
    \item understand subjective probabilities;
    \item define utility functions;
    \item distinguish risk-neutral, risk-averse, and risk-seeking behavior;
    \item define expected utility;
    \item use certainty equivalents;
    \item define risk premiums;
    \item recognize diminishing marginal utility;
    \item understand stochastic dominance conceptually;
    \item distinguish first-order and second-order stochastic dominance;
    \item understand Bayesian decision theory;
    \item update beliefs using Bayes' theorem;
    \item construct decision trees;
    \item evaluate decision trees using rollback analysis;
    \item distinguish decision nodes and chance nodes;
    \item calculate the expected value of perfect information;
    \item calculate the expected value of sample information;
    \item understand the value of imperfect information;
    \item understand prior and posterior probabilities;
    \item formulate Bayes-risk minimization;
    \item understand loss functions in statistical decisions;
    \item recognize classification as a decision problem;
    \item understand threshold decision rules;
    \item define sensitivity and specificity conceptually;
    \item understand false-positive and false-negative costs;
    \item recognize receiver operating characteristic tradeoffs
          conceptually;
    \item understand utility under multiple attributes;
    \item recognize multiattribute utility theory;
    \item distinguish decision theory from game theory;
    \item recognize sequential decision problems;
    \item connect decision theory with dynamic programming;
    \item understand ambiguity and robust decision making conceptually;
    \item understand behavioral departures from classical expected
          utility conceptually;
    \item connect decision theory with economics, medicine, engineering,
          management, finance, public policy, and operations research.
\end{itemize}
}

%%%%%%%%%%%%%%%%%%%%%%%%%%%%%%%%%%%%%%%%%%%%%%%%%%%%%%%%%%%%
\subsection{What Is Decision Theory?}
\label{subsec:what-is-decision-theory}
%%%%%%%%%%%%%%%%%%%%%%%%%%%%%%%%%%%%%%%%%%%%%%%%%%%%%%%%%%%%

Decision theory studies how choices should be made when outcomes depend on
both:

\[
\boxed{
\text{the decision selected}
}
\]

and

\[
\boxed{
\text{uncertain external conditions}.
}
\]

A decision problem contains:

\[
\boxed{
\text{alternatives}
+
\text{states of nature}
+
\text{consequences}
+
\text{preferences}.
}
\]

%%%%%%%%%%%%%%%%%%%%%%%%%%%%%%%%%%%%%%%%%%%%%%%%%%%%%%%%%%%%
\subsection{Decision Alternatives}
\label{subsec:decision-alternatives}
%%%%%%%%%%%%%%%%%%%%%%%%%%%%%%%%%%%%%%%%%%%%%%%%%%%%%%%%%%%%

Let

\[
\boxed{
A
=
\{a_1,a_2,\ldots,a_m\}
}
\]

be the set of available actions.

Examples include:

\[
\boxed{
\text{build or not build},
}
\]

\[
\boxed{
\text{order different quantities},
}
\]

\[
\boxed{
\text{choose an investment},
}
\]

or

\[
\boxed{
\text{select a medical treatment}.
}
\]

%%%%%%%%%%%%%%%%%%%%%%%%%%%%%%%%%%%%%%%%%%%%%%%%%%%%%%%%%%%%
\subsection{States of Nature}
\label{subsec:states-of-nature}
%%%%%%%%%%%%%%%%%%%%%%%%%%%%%%%%%%%%%%%%%%%%%%%%%%%%%%%%%%%%

A state of nature is an uncertain condition outside the decision maker's
direct control.

Let

\[
\boxed{
\Theta
=
\{\theta_1,\theta_2,\ldots,\theta_n\}.
}
\]

Examples include:

\[
\boxed{
\text{high or low demand},
}
\]

\[
\boxed{
\text{market growth or decline},
}
\]

or

\[
\boxed{
\text{disease present or absent}.
}
\]

%%%%%%%%%%%%%%%%%%%%%%%%%%%%%%%%%%%%%%%%%%%%%%%%%%%%%%%%%%%%
\subsection{Consequences}
\label{subsec:decision-consequences}
%%%%%%%%%%%%%%%%%%%%%%%%%%%%%%%%%%%%%%%%%%%%%%%%%%%%%%%%%%%%

Each pair

\[
(a_i,\theta_j)
\]

produces a consequence.

A numerical payoff may be written

\[
\boxed{
u_{ij}
=
u(a_i,\theta_j).
}
\]

Alternatively, a loss may be written

\[
\boxed{
L(a_i,\theta_j).
}
\]

The decision criterion determines how these consequences are compared.

%%%%%%%%%%%%%%%%%%%%%%%%%%%%%%%%%%%%%%%%%%%%%%%%%%%%%%%%%%%%
\subsection{Payoff Table}
\label{subsec:payoff-table}
%%%%%%%%%%%%%%%%%%%%%%%%%%%%%%%%%%%%%%%%%%%%%%%%%%%%%%%%%%%%

A decision problem may be summarized by a payoff table:

\[
\begin{array}{c|ccc}
&
\theta_1
&
\theta_2
&
\theta_3
\\
\hline
a_1
&
u_{11}
&
u_{12}
&
u_{13}
\\
a_2
&
u_{21}
&
u_{22}
&
u_{23}
\\
a_3
&
u_{31}
&
u_{32}
&
u_{33}
\end{array}
\]

Rows correspond to decisions.

Columns correspond to states of nature.

%%%%%%%%%%%%%%%%%%%%%%%%%%%%%%%%%%%%%%%%%%%%%%%%%%%%%%%%%%%%
\subsection{Loss Table}
\label{subsec:loss-table}
%%%%%%%%%%%%%%%%%%%%%%%%%%%%%%%%%%%%%%%%%%%%%%%%%%%%%%%%%%%%

If smaller values are preferred, use a loss table:

\[
\boxed{
L_{ij}
=
L(a_i,\theta_j).
}
\]

The optimization goal is then typically:

\[
\boxed{
\text{minimize loss}.
}
\]

Payoff and loss formulations are equivalent when one is transformed
consistently into the other.

%%%%%%%%%%%%%%%%%%%%%%%%%%%%%%%%%%%%%%%%%%%%%%%%%%%%%%%%%%%%
\subsection{Decision Making Under Certainty}
\label{subsec:decision-under-certainty}
%%%%%%%%%%%%%%%%%%%%%%%%%%%%%%%%%%%%%%%%%%%%%%%%%%%%%%%%%%%%

Under certainty, the state of nature is known.

If state

\[
\theta_j
\]

is known, simply choose

\[
\boxed{
a^*
\in
\operatorname*{arg\,max}_{a_i}
u(a_i,\theta_j)
}
\]

for a payoff problem.

No probability model is needed.

%%%%%%%%%%%%%%%%%%%%%%%%%%%%%%%%%%%%%%%%%%%%%%%%%%%%%%%%%%%%
\subsection{Decision Making Under Risk}
\label{subsec:decision-under-risk}
%%%%%%%%%%%%%%%%%%%%%%%%%%%%%%%%%%%%%%%%%%%%%%%%%%%%%%%%%%%%

Under risk, the states of nature are uncertain but probabilities are
available.

Suppose

\[
\boxed{
\mathbb{P}(\theta_j)
=
p_j,
}
\]

with

\[
p_j\geq0,
\]

and

\[
\sum_jp_j=1.
\]

Then probability-based criteria such as expected value or expected utility
can be used.

%%%%%%%%%%%%%%%%%%%%%%%%%%%%%%%%%%%%%%%%%%%%%%%%%%%%%%%%%%%%
\subsection{Decision Making Under Uncertainty}
\label{subsec:decision-under-uncertainty}
%%%%%%%%%%%%%%%%%%%%%%%%%%%%%%%%%%%%%%%%%%%%%%%%%%%%%%%%%%%%

Under uncertainty, meaningful probabilities may not be available.

One may then use criteria such as:

\[
\boxed{
\text{maximin},
}
\]

\[
\boxed{
\text{maximax},
}
\]

\[
\boxed{
\text{minimax regret},
}
\]

\[
\boxed{
\text{Hurwicz},
}
\]

or

\[
\boxed{
\text{Laplace}.
}
\]

Different criteria represent different attitudes toward uncertainty.

%%%%%%%%%%%%%%%%%%%%%%%%%%%%%%%%%%%%%%%%%%%%%%%%%%%%%%%%%%%%
\subsection{Expected Monetary Value}
\label{subsec:expected-monetary-value}
%%%%%%%%%%%%%%%%%%%%%%%%%%%%%%%%%%%%%%%%%%%%%%%%%%%%%%%%%%%%

Suppose action \(a_i\) produces monetary payoff \(u_{ij}\) in state
\(\theta_j\).

Its expected monetary value is

\[
\boxed{
EMV(a_i)
=
\sum_{j=1}^{n}
p_j
u_{ij}.
}
\]

A risk-neutral decision maker chooses an action maximizing EMV.

%%%%%%%%%%%%%%%%%%%%%%%%%%%%%%%%%%%%%%%%%%%%%%%%%%%%%%%%%%%%
\subsection{Worked Example: Expected Monetary Value}
\label{subsec:worked-emv}
%%%%%%%%%%%%%%%%%%%%%%%%%%%%%%%%%%%%%%%%%%%%%%%%%%%%%%%%%%%%

\begin{workedexamplebox}{Choosing a Production Plan}

Suppose demand may be high or low.

The probability of high demand is

\[
0.6,
\]

and low demand is

\[
0.4.
\]

Plan \(A\) has profits

\[
100
\]

under high demand and

\[
40
\]

under low demand.

Plan \(B\) has profits

\[
80
\]

under high demand and

\[
60
\]

under low demand.

Then

\[
EMV(A)
=
0.6(100)
+
0.4(40)
=
76.
\]

For plan \(B\),

\[
EMV(B)
=
0.6(80)
+
0.4(60)
=
72.
\]

Therefore,

\begin{answerbox}
\[
\boxed{
\text{Plan }A
}
\]

has the larger expected monetary value.
\end{answerbox}

\end{workedexamplebox}

%%%%%%%%%%%%%%%%%%%%%%%%%%%%%%%%%%%%%%%%%%%%%%%%%%%%%%%%%%%%
\subsection{Expected Loss}
\label{subsec:expected-loss}
%%%%%%%%%%%%%%%%%%%%%%%%%%%%%%%%%%%%%%%%%%%%%%%%%%%%%%%%%%%%

If \(L_{ij}\) is a loss,

\[
\boxed{
\mathbb{E}
[
L(a_i,\Theta)
]
=
\sum_j
p_jL_{ij}.
}
\]

The Bayes action for a known probability model minimizes expected loss:

\[
\boxed{
a^*
\in
\operatorname*{arg\,min}_a
\mathbb{E}
[
L(a,\Theta)
].
}
\]

%%%%%%%%%%%%%%%%%%%%%%%%%%%%%%%%%%%%%%%%%%%%%%%%%%%%%%%%%%%%
\subsection{Risk-Neutral Decision Making}
\label{subsec:risk-neutral-decision-making}
%%%%%%%%%%%%%%%%%%%%%%%%%%%%%%%%%%%%%%%%%%%%%%%%%%%%%%%%%%%%

A risk-neutral decision maker evaluates monetary alternatives according to
expected monetary value.

Thus two uncertain payoffs with the same expectation are treated as
equally desirable if all other considerations are ignored.

This does not describe every real decision maker.

%%%%%%%%%%%%%%%%%%%%%%%%%%%%%%%%%%%%%%%%%%%%%%%%%%%%%%%%%%%%
\subsection{Maximax Criterion}
\label{subsec:maximax-criterion}
%%%%%%%%%%%%%%%%%%%%%%%%%%%%%%%%%%%%%%%%%%%%%%%%%%%%%%%%%%%%

For each action, identify its best possible payoff:

\[
\boxed{
M_i
=
\max_j
u_{ij}.
}
\]

The maximax rule chooses

\[
\boxed{
a^*
\in
\operatorname*{arg\,max}_i
M_i.
}
\]

This criterion emphasizes the best possible outcome.

%%%%%%%%%%%%%%%%%%%%%%%%%%%%%%%%%%%%%%%%%%%%%%%%%%%%%%%%%%%%
\subsection{Maximin Criterion}
\label{subsec:maximin-criterion}
%%%%%%%%%%%%%%%%%%%%%%%%%%%%%%%%%%%%%%%%%%%%%%%%%%%%%%%%%%%%

For each action, identify its worst payoff:

\[
\boxed{
m_i
=
\min_j
u_{ij}.
}
\]

The maximin rule chooses

\[
\boxed{
a^*
\in
\operatorname*{arg\,max}_i
m_i.
}
\]

This is a conservative payoff criterion.

%%%%%%%%%%%%%%%%%%%%%%%%%%%%%%%%%%%%%%%%%%%%%%%%%%%%%%%%%%%%
\subsection{Worked Example: Maximin and Maximax}
\label{subsec:worked-maximin-maximax}
%%%%%%%%%%%%%%%%%%%%%%%%%%%%%%%%%%%%%%%%%%%%%%%%%%%%%%%%%%%%

\begin{workedexamplebox}{Two Uncertainty Criteria}

Suppose:

\[
\begin{array}{c|ccc}
&
\theta_1
&
\theta_2
&
\theta_3
\\
\hline
A
&
40
&
60
&
90
\\
B
&
20
&
80
&
120
\\
C
&
50
&
55
&
65
\end{array}
\]

The minimum payoff for each action is

\[
A:40,
\qquad
B:20,
\qquad
C:50.
\]

Therefore maximin selects

\[
\boxed{
C.
}
\]

The maximum payoff for each action is

\[
A:90,
\qquad
B:120,
\qquad
C:65.
\]

Therefore maximax selects

\[
\boxed{
B.
}
\]

\begin{answerbox}
\[
\boxed{
\text{Maximin: }C,
\qquad
\text{Maximax: }B.
}
\]
\end{answerbox}

\end{workedexamplebox}

%%%%%%%%%%%%%%%%%%%%%%%%%%%%%%%%%%%%%%%%%%%%%%%%%%%%%%%%%%%%
\subsection{Hurwicz Criterion}
\label{subsec:hurwicz-criterion}
%%%%%%%%%%%%%%%%%%%%%%%%%%%%%%%%%%%%%%%%%%%%%%%%%%%%%%%%%%%%

The Hurwicz criterion combines best-case and worst-case outcomes.

Let

\[
\alpha\in[0,1]
\]

be an optimism coefficient.

For a payoff problem,

\[
\boxed{
H_i
=
\alpha
\max_j
u_{ij}
+
(1-\alpha)
\min_j
u_{ij}.
}
\]

Choose the action with largest \(H_i\).

%%%%%%%%%%%%%%%%%%%%%%%%%%%%%%%%%%%%%%%%%%%%%%%%%%%%%%%%%%%%
\subsection{Interpretation of Hurwicz Parameter}
\label{subsec:hurwicz-parameter}
%%%%%%%%%%%%%%%%%%%%%%%%%%%%%%%%%%%%%%%%%%%%%%%%%%%%%%%%%%%%

If

\[
\alpha=1,
\]

Hurwicz reduces to maximax.

If

\[
\alpha=0,
\]

it reduces to maximin.

Intermediate values represent a compromise between optimistic and
pessimistic evaluation.

%%%%%%%%%%%%%%%%%%%%%%%%%%%%%%%%%%%%%%%%%%%%%%%%%%%%%%%%%%%%
\subsection{Laplace Principle}
\label{subsec:laplace-principle}
%%%%%%%%%%%%%%%%%%%%%%%%%%%%%%%%%%%%%%%%%%%%%%%%%%%%%%%%%%%%

If no probabilities are known, the Laplace principle assigns equal
probability to each state:

\[
\boxed{
p_j
=
\frac1n.
}
\]

Then choose the action maximizing average payoff:

\[
\boxed{
\frac1n
\sum_j
u_{ij}.
}
\]

This is a modeling assumption, not a consequence of ignorance alone.

%%%%%%%%%%%%%%%%%%%%%%%%%%%%%%%%%%%%%%%%%%%%%%%%%%%%%%%%%%%%
\subsection{Regret}
\label{subsec:decision-regret}
%%%%%%%%%%%%%%%%%%%%%%%%%%%%%%%%%%%%%%%%%%%%%%%%%%%%%%%%%%%%

Regret measures the opportunity lost because a decision was not the best
one for the realized state.

For a payoff problem, let

\[
u_j^*
=
\max_i
u_{ij}.
\]

Then regret is

\[
\boxed{
R_{ij}
=
u_j^*
-
u_{ij}.
}
\]

Regret is always nonnegative.

%%%%%%%%%%%%%%%%%%%%%%%%%%%%%%%%%%%%%%%%%%%%%%%%%%%%%%%%%%%%
\subsection{Regret Table}
\label{subsec:regret-table}
%%%%%%%%%%%%%%%%%%%%%%%%%%%%%%%%%%%%%%%%%%%%%%%%%%%%%%%%%%%%

A regret table is obtained by subtracting each payoff in a column from the
best payoff in that state.

If

\[
\max_i
u_{ij}
=
100,
\]

and action \(a_i\) produces

\[
70,
\]

then its regret is

\[
\boxed{
30.
}
\]

%%%%%%%%%%%%%%%%%%%%%%%%%%%%%%%%%%%%%%%%%%%%%%%%%%%%%%%%%%%%
\subsection{Minimax Regret}
\label{subsec:minimax-regret}
%%%%%%%%%%%%%%%%%%%%%%%%%%%%%%%%%%%%%%%%%%%%%%%%%%%%%%%%%%%%

For each action, calculate its maximum regret:

\[
\boxed{
R_i^{\max}
=
\max_j
R_{ij}.
}
\]

Then choose

\[
\boxed{
a^*
\in
\operatorname*{arg\,min}_i
R_i^{\max}.
}
\]

This minimizes the worst opportunity loss.

%%%%%%%%%%%%%%%%%%%%%%%%%%%%%%%%%%%%%%%%%%%%%%%%%%%%%%%%%%%%
\subsection{Worked Example: Minimax Regret}
\label{subsec:worked-minimax-regret}
%%%%%%%%%%%%%%%%%%%%%%%%%%%%%%%%%%%%%%%%%%%%%%%%%%%%%%%%%%%%

\begin{workedexamplebox}{Constructing a Regret Table}

Suppose the payoff table is

\[
\begin{array}{c|cc}
&
\theta_1
&
\theta_2
\\
\hline
A
&
90
&
40
\\
B
&
70
&
70
\\
C
&
50
&
100
\end{array}
\]

The best payoff in state \(\theta_1\) is

\[
90,
\]

and in state \(\theta_2\) it is

\[
100.
\]

Thus regrets are

\[
\begin{array}{c|cc}
&
\theta_1
&
\theta_2
\\
\hline
A
&
0
&
60
\\
B
&
20
&
30
\\
C
&
40
&
0
\end{array}
\]

Maximum regrets are

\[
A:60,
\qquad
B:30,
\qquad
C:40.
\]

Therefore,

\begin{answerbox}
\[
\boxed{
B
}
\]

is the minimax-regret decision.
\end{answerbox}

\end{workedexamplebox}

%%%%%%%%%%%%%%%%%%%%%%%%%%%%%%%%%%%%%%%%%%%%%%%%%%%%%%%%%%%%
\subsection{Utility}
\label{subsec:decision-utility}
%%%%%%%%%%%%%%%%%%%%%%%%%%%%%%%%%%%%%%%%%%%%%%%%%%%%%%%%%%%%

A utility function assigns a numerical preference value to outcomes.

If wealth is \(w\),

\[
\boxed{
U(w)
}
\]

represents the decision maker's preference for that wealth level.

Utility values are meaningful through preference comparisons and expected
utility calculations rather than as physical units.

%%%%%%%%%%%%%%%%%%%%%%%%%%%%%%%%%%%%%%%%%%%%%%%%%%%%%%%%%%%%
\subsection{Expected Utility}
\label{subsec:expected-utility}
%%%%%%%%%%%%%%%%%%%%%%%%%%%%%%%%%%%%%%%%%%%%%%%%%%%%%%%%%%%%

Suppose outcome \(x_j\) occurs with probability \(p_j\).

Expected utility is

\[
\boxed{
EU
=
\sum_j
p_j
U(x_j).
}
\]

An expected-utility decision maker selects the action with the highest
expected utility.

%%%%%%%%%%%%%%%%%%%%%%%%%%%%%%%%%%%%%%%%%%%%%%%%%%%%%%%%%%%%
\subsection{Expected Utility Versus Utility of Expected Wealth}
\label{subsec:expected-utility-versus-expected-wealth}
%%%%%%%%%%%%%%%%%%%%%%%%%%%%%%%%%%%%%%%%%%%%%%%%%%%%%%%%%%%%

In general,

\[
\boxed{
\mathbb{E}
[
U(W)
]
\neq
U
(
\mathbb{E}[W]
).
}
\]

The difference is essential to modeling risk preferences.

Only when utility is affine do expected utility and expected monetary value
produce the same ordering of lotteries.

%%%%%%%%%%%%%%%%%%%%%%%%%%%%%%%%%%%%%%%%%%%%%%%%%%%%%%%%%%%%
\subsection{Risk Aversion}
\label{subsec:risk-aversion}
%%%%%%%%%%%%%%%%%%%%%%%%%%%%%%%%%%%%%%%%%%%%%%%%%%%%%%%%%%%%

A decision maker is risk averse if they prefer a certain outcome to a fair
risky alternative with the same expected monetary value.

A differentiable utility function is risk averse when it is concave:

\[
\boxed{
U''(w)<0.
}
\]

Concavity represents diminishing marginal utility of wealth.

%%%%%%%%%%%%%%%%%%%%%%%%%%%%%%%%%%%%%%%%%%%%%%%%%%%%%%%%%%%%
\subsection{Risk Neutrality}
\label{subsec:risk-neutrality}
%%%%%%%%%%%%%%%%%%%%%%%%%%%%%%%%%%%%%%%%%%%%%%%%%%%%%%%%%%%%

Risk neutrality corresponds to affine utility:

\[
\boxed{
U(w)
=
aw+b,
\qquad
a>0.
}
\]

Then maximizing expected utility is equivalent to maximizing expected
wealth.

%%%%%%%%%%%%%%%%%%%%%%%%%%%%%%%%%%%%%%%%%%%%%%%%%%%%%%%%%%%%
\subsection{Risk Seeking}
\label{subsec:risk-seeking}
%%%%%%%%%%%%%%%%%%%%%%%%%%%%%%%%%%%%%%%%%%%%%%%%%%%%%%%%%%%%

A risk-seeking decision maker may prefer a fair risky alternative to the
certain expected value.

This behavior can be represented by convex utility:

\[
\boxed{
U''(w)>0.
}
\]

Real preferences may vary across decision domains.

%%%%%%%%%%%%%%%%%%%%%%%%%%%%%%%%%%%%%%%%%%%%%%%%%%%%%%%%%%%%
\subsection{Jensen's Inequality and Risk Aversion}
\label{subsec:jensen-risk-aversion}
%%%%%%%%%%%%%%%%%%%%%%%%%%%%%%%%%%%%%%%%%%%%%%%%%%%%%%%%%%%%

For a concave utility function,

\[
\boxed{
\mathbb{E}
[
U(W)
]
\leq
U
(
\mathbb{E}[W]
).
}
\]

Thus the utility of expected wealth exceeds or equals expected utility.

This inequality mathematically expresses aversion to risk.

%%%%%%%%%%%%%%%%%%%%%%%%%%%%%%%%%%%%%%%%%%%%%%%%%%%%%%%%%%%%
\subsection{Certainty Equivalent}
\label{subsec:certainty-equivalent}
%%%%%%%%%%%%%%%%%%%%%%%%%%%%%%%%%%%%%%%%%%%%%%%%%%%%%%%%%%%%

The certainty equivalent \(CE\) of a risky payoff \(W\) satisfies

\[
\boxed{
U(CE)
=
\mathbb{E}
[
U(W)
].
}
\]

Therefore,

\[
\boxed{
CE
=
U^{-1}
\left(
\mathbb{E}[U(W)]
\right).
}
\]

It is the certain amount the decision maker considers equivalent to the
risky prospect.

%%%%%%%%%%%%%%%%%%%%%%%%%%%%%%%%%%%%%%%%%%%%%%%%%%%%%%%%%%%%
\subsection{Risk Premium}
\label{subsec:risk-premium}
%%%%%%%%%%%%%%%%%%%%%%%%%%%%%%%%%%%%%%%%%%%%%%%%%%%%%%%%%%%%

The risk premium is

\[
\boxed{
\pi
=
\mathbb{E}[W]
-
CE.
}
\]

For a risk-averse decision maker,

\[
\boxed{
\pi\geq0
}
\]

under standard conditions.

It represents the maximum expected wealth sacrificed to eliminate the
risk.

%%%%%%%%%%%%%%%%%%%%%%%%%%%%%%%%%%%%%%%%%%%%%%%%%%%%%%%%%%%%
\subsection{Worked Example: Certainty Equivalent}
\label{subsec:worked-certainty-equivalent}
%%%%%%%%%%%%%%%%%%%%%%%%%%%%%%%%%%%%%%%%%%%%%%%%%%%%%%%%%%%%

\begin{workedexamplebox}{Square-Root Utility}

Suppose

\[
U(w)=\sqrt{w}.
\]

A lottery gives

\[
100
\]

with probability

\[
0.5
\]

and

\[
400
\]

with probability

\[
0.5.
\]

Expected utility is

\[
EU
=
0.5(10)
+
0.5(20)
=
15.
\]

The certainty equivalent satisfies

\[
\sqrt{CE}
=
15.
\]

Thus,

\[
CE
=
225.
\]

Expected wealth is

\[
\mathbb{E}[W]
=
0.5(100)
+
0.5(400)
=
250.
\]

Therefore the risk premium is

\[
250-225
=
25.
\]

\begin{answerbox}
\[
\boxed{
CE=225,
\qquad
\pi=25.
}
\]
\end{answerbox}

\end{workedexamplebox}

%%%%%%%%%%%%%%%%%%%%%%%%%%%%%%%%%%%%%%%%%%%%%%%%%%%%%%%%%%%%
\subsection{Absolute Risk Aversion}
\label{subsec:absolute-risk-aversion}
%%%%%%%%%%%%%%%%%%%%%%%%%%%%%%%%%%%%%%%%%%%%%%%%%%%%%%%%%%%%

For a twice-differentiable utility function with \(U'(w)>0\), the
Arrow--Pratt coefficient of absolute risk aversion is

\[
\boxed{
A(w)
=
-
\frac{
U''(w)
}{
U'(w)
}.
}
\]

Larger values indicate greater local aversion to small absolute risks.

%%%%%%%%%%%%%%%%%%%%%%%%%%%%%%%%%%%%%%%%%%%%%%%%%%%%%%%%%%%%
\subsection{Relative Risk Aversion}
\label{subsec:relative-risk-aversion}
%%%%%%%%%%%%%%%%%%%%%%%%%%%%%%%%%%%%%%%%%%%%%%%%%%%%%%%%%%%%

Relative risk aversion is

\[
\boxed{
R(w)
=
-
w
\frac{
U''(w)
}{
U'(w)
}.
}
\]

It measures sensitivity to risks expressed relative to wealth.

These measures are widely used in economic and financial decision theory.

%%%%%%%%%%%%%%%%%%%%%%%%%%%%%%%%%%%%%%%%%%%%%%%%%%%%%%%%%%%%
\subsection{Common Utility Functions}
\label{subsec:common-utility-functions}
%%%%%%%%%%%%%%%%%%%%%%%%%%%%%%%%%%%%%%%%%%%%%%%%%%%%%%%%%%%%

Examples include logarithmic utility:

\[
\boxed{
U(w)=\ln w,
\qquad
w>0,
}
\]

power utility:

\[
\boxed{
U(w)
=
\frac{
w^{1-\gamma}
}{
1-\gamma
},
}
\]

for appropriate \(\gamma\), and exponential utility:

\[
\boxed{
U(w)
=
-e^{-aw},
\qquad
a>0.
}
\]

Each expresses a different risk-preference structure.

%%%%%%%%%%%%%%%%%%%%%%%%%%%%%%%%%%%%%%%%%%%%%%%%%%%%%%%%%%%%
\subsection{Expected Utility Axioms Preview}
\label{subsec:expected-utility-axioms}
%%%%%%%%%%%%%%%%%%%%%%%%%%%%%%%%%%%%%%%%%%%%%%%%%%%%%%%%%%%%

Classical expected-utility theory can be derived from preference axioms
over lotteries.

These include concepts such as:

\[
\boxed{
\text{completeness},
}
\]

\[
\boxed{
\text{transitivity},
}
\]

\[
\boxed{
\text{continuity},
}
\]

and

\[
\boxed{
\text{independence}.
}
\]

Under suitable axioms, preferences can be represented by expected utility.

%%%%%%%%%%%%%%%%%%%%%%%%%%%%%%%%%%%%%%%%%%%%%%%%%%%%%%%%%%%%
\subsection{Utility Transformations}
\label{subsec:utility-transformations}
%%%%%%%%%%%%%%%%%%%%%%%%%%%%%%%%%%%%%%%%%%%%%%%%%%%%%%%%%%%%

If \(U\) represents expected-utility preferences, then

\[
\boxed{
\widetilde{U}(x)
=
aU(x)+b,
\qquad
a>0,
}
\]

represents the same preferences.

Thus cardinal utility in expected-utility theory is unique only up to
positive affine transformation.

%%%%%%%%%%%%%%%%%%%%%%%%%%%%%%%%%%%%%%%%%%%%%%%%%%%%%%%%%%%%
\subsection{Stochastic Dominance}
\label{subsec:stochastic-dominance}
%%%%%%%%%%%%%%%%%%%%%%%%%%%%%%%%%%%%%%%%%%%%%%%%%%%%%%%%%%%%

Stochastic dominance compares random outcomes without specifying one exact
utility function.

It asks whether one distribution is preferred by an entire class of
decision makers.

%%%%%%%%%%%%%%%%%%%%%%%%%%%%%%%%%%%%%%%%%%%%%%%%%%%%%%%%%%%%
\subsection{First-Order Stochastic Dominance}
\label{subsec:first-order-stochastic-dominance}
%%%%%%%%%%%%%%%%%%%%%%%%%%%%%%%%%%%%%%%%%%%%%%%%%%%%%%%%%%%%

Random variable \(X\) first-order stochastically dominates \(Y\) if

\[
\boxed{
F_X(t)
\leq
F_Y(t)
}
\]

for all \(t\), with strict inequality somewhere.

Then every increasing utility function prefers \(X\) to \(Y\).

Intuitively, \(X\) tends to produce larger outcomes.

%%%%%%%%%%%%%%%%%%%%%%%%%%%%%%%%%%%%%%%%%%%%%%%%%%%%%%%%%%%%
\subsection{Second-Order Stochastic Dominance}
\label{subsec:second-order-stochastic-dominance}
%%%%%%%%%%%%%%%%%%%%%%%%%%%%%%%%%%%%%%%%%%%%%%%%%%%%%%%%%%%%

Second-order stochastic dominance is relevant to increasing concave utility
functions.

One common condition is

\[
\boxed{
\int_{-\infty}^{t}
F_X(s)\,ds
\leq
\int_{-\infty}^{t}
F_Y(s)\,ds
}
\]

for every \(t\), with suitable endpoint conditions.

It captures preference among risk-averse expected-utility decision makers.

%%%%%%%%%%%%%%%%%%%%%%%%%%%%%%%%%%%%%%%%%%%%%%%%%%%%%%%%%%%%
\subsection{Decision Trees}
\label{subsec:decision-trees}
%%%%%%%%%%%%%%%%%%%%%%%%%%%%%%%%%%%%%%%%%%%%%%%%%%%%%%%%%%%%

A decision tree represents sequential decisions and uncertain events.

Typical node types are:

\[
\boxed{
\text{decision nodes}
}
\]

and

\[
\boxed{
\text{chance nodes}.
}
\]

Branches represent available choices or possible outcomes.

%%%%%%%%%%%%%%%%%%%%%%%%%%%%%%%%%%%%%%%%%%%%%%%%%%%%%%%%%%%%
\subsection{Decision Nodes}
\label{subsec:decision-nodes}
%%%%%%%%%%%%%%%%%%%%%%%%%%%%%%%%%%%%%%%%%%%%%%%%%%%%%%%%%%%%

At a decision node, the decision maker chooses one branch.

The value of a payoff decision node is

\[
\boxed{
\max
\{
\text{values of available branches}
\}.
}
\]

For a cost problem, use the minimum instead.

%%%%%%%%%%%%%%%%%%%%%%%%%%%%%%%%%%%%%%%%%%%%%%%%%%%%%%%%%%%%
\subsection{Chance Nodes}
\label{subsec:chance-nodes}
%%%%%%%%%%%%%%%%%%%%%%%%%%%%%%%%%%%%%%%%%%%%%%%%%%%%%%%%%%%%

At a chance node, branches occur randomly.

Its value is the probability-weighted average:

\[
\boxed{
\sum_j
p_jV_j.
}
\]

Thus decision trees alternate:

\[
\boxed{
\text{optimization}
}
\]

and

\[
\boxed{
\text{expectation}.
}
\]

%%%%%%%%%%%%%%%%%%%%%%%%%%%%%%%%%%%%%%%%%%%%%%%%%%%%%%%%%%%%
\subsection{Rollback Analysis}
\label{subsec:rollback-analysis}
%%%%%%%%%%%%%%%%%%%%%%%%%%%%%%%%%%%%%%%%%%%%%%%%%%%%%%%%%%%%

Decision trees are evaluated from right to left.

This is called rollback or fold-back analysis.

At each chance node:

\[
\boxed{
\text{take expectation}.
}
\]

At each decision node:

\[
\boxed{
\text{choose the best branch}.
}
\]

This is closely related to backward induction in dynamic programming.

%%%%%%%%%%%%%%%%%%%%%%%%%%%%%%%%%%%%%%%%%%%%%%%%%%%%%%%%%%%%
\subsection{Worked Example: Decision Tree}
\label{subsec:worked-decision-tree}
%%%%%%%%%%%%%%%%%%%%%%%%%%%%%%%%%%%%%%%%%%%%%%%%%%%%%%%%%%%%

\begin{workedexamplebox}{Investment Decision}

Suppose investing produces:

\[
100
\]

with probability

\[
0.7,
\]

and

\[
-40
\]

with probability

\[
0.3.
\]

Not investing produces

\[
20
\]

with certainty.

The expected value of investing is

\[
0.7(100)
+
0.3(-40)
=
58.
\]

Since

\[
58>20,
\]

a risk-neutral rollback selects

\begin{answerbox}
\[
\boxed{
\text{Invest}.
}
\]
\end{answerbox}

\end{workedexamplebox}

%%%%%%%%%%%%%%%%%%%%%%%%%%%%%%%%%%%%%%%%%%%%%%%%%%%%%%%%%%%%
\subsection{Perfect Information}
\label{subsec:perfect-information-decision}
%%%%%%%%%%%%%%%%%%%%%%%%%%%%%%%%%%%%%%%%%%%%%%%%%%%%%%%%%%%%

Perfect information reveals the true state of nature before the decision
is made.

The decision maker can then choose the best action separately for each
state.

For payoff maximization,

\[
\boxed{
EV_{\mathrm{PI}}
=
\sum_j
p_j
\max_i
u_{ij}.
}
\]

%%%%%%%%%%%%%%%%%%%%%%%%%%%%%%%%%%%%%%%%%%%%%%%%%%%%%%%%%%%%
\subsection{Expected Value Without Perfect Information}
\label{subsec:expected-value-without-pi}
%%%%%%%%%%%%%%%%%%%%%%%%%%%%%%%%%%%%%%%%%%%%%%%%%%%%%%%%%%%%

Without additional information, choose the action with highest expected
value:

\[
\boxed{
EV_{\mathrm{current}}
=
\max_i
\sum_j
p_j
u_{ij}.
}
\]

Perfect information cannot reduce the optimal expected payoff.

%%%%%%%%%%%%%%%%%%%%%%%%%%%%%%%%%%%%%%%%%%%%%%%%%%%%%%%%%%%%
\subsection{Expected Value of Perfect Information}
\label{subsec:evpi-decision-theory}
%%%%%%%%%%%%%%%%%%%%%%%%%%%%%%%%%%%%%%%%%%%%%%%%%%%%%%%%%%%%

For payoff maximization,

\[
\boxed{
EVPI
=
EV_{\mathrm{PI}}
-
EV_{\mathrm{current}}.
}
\]

Thus,

\[
\boxed{
EVPI\geq0.
}
\]

EVPI is the maximum expected amount worth paying for perfect information.

%%%%%%%%%%%%%%%%%%%%%%%%%%%%%%%%%%%%%%%%%%%%%%%%%%%%%%%%%%%%
\subsection{Worked Example: EVPI}
\label{subsec:worked-evpi-decision}
%%%%%%%%%%%%%%%%%%%%%%%%%%%%%%%%%%%%%%%%%%%%%%%%%%%%%%%%%%%%

\begin{workedexamplebox}{Value of Knowing Demand}

Suppose:

\[
\begin{array}{c|cc}
&
\text{High}
&
\text{Low}
\\
\hline
A
&
100
&
20
\\
B
&
70
&
60
\end{array}
\]

with

\[
P(\text{High})=0.6,
\qquad
P(\text{Low})=0.4.
\]

Without perfect information:

\[
EMV(A)
=
0.6(100)+0.4(20)
=
68,
\]

\[
EMV(B)
=
0.6(70)+0.4(60)
=
66.
\]

Thus,

\[
EV_{\mathrm{current}}
=
68.
\]

With perfect information, choose \(A\) under high demand and \(B\) under
low demand:

\[
EV_{\mathrm{PI}}
=
0.6(100)
+
0.4(60)
=
84.
\]

Therefore,

\begin{answerbox}
\[
\boxed{
EVPI
=
84-68
=
16.
}
\]
\end{answerbox}

\end{workedexamplebox}

%%%%%%%%%%%%%%%%%%%%%%%%%%%%%%%%%%%%%%%%%%%%%%%%%%%%%%%%%%%%
\subsection{Sample Information}
\label{subsec:sample-information}
%%%%%%%%%%%%%%%%%%%%%%%%%%%%%%%%%%%%%%%%%%%%%%%%%%%%%%%%%%%%

Often information is imperfect.

Examples include:

\[
\boxed{
\text{market surveys},
}
\]

\[
\boxed{
\text{medical tests},
}
\]

\[
\boxed{
\text{pilot studies},
}
\]

or

\[
\boxed{
\text{weather forecasts}.
}
\]

The information changes probabilities rather than revealing the state
with certainty.

%%%%%%%%%%%%%%%%%%%%%%%%%%%%%%%%%%%%%%%%%%%%%%%%%%%%%%%%%%%%
\subsection{Prior Probabilities}
\label{subsec:prior-probabilities}
%%%%%%%%%%%%%%%%%%%%%%%%%%%%%%%%%%%%%%%%%%%%%%%%%%%%%%%%%%%%

Before observing new information, uncertainty is described by prior
probabilities:

\[
\boxed{
P(\theta_j).
}
\]

These may come from:

\[
\boxed{
\text{historical data},
}
\]

\[
\boxed{
\text{expert judgment},
}
\]

or

\[
\boxed{
\text{previous analysis}.
}
\]

%%%%%%%%%%%%%%%%%%%%%%%%%%%%%%%%%%%%%%%%%%%%%%%%%%%%%%%%%%%%
\subsection{Conditional Information Model}
\label{subsec:conditional-information-model}
%%%%%%%%%%%%%%%%%%%%%%%%%%%%%%%%%%%%%%%%%%%%%%%%%%%%%%%%%%%%

Let \(Y\) denote an observed signal.

The information model is specified by probabilities such as

\[
\boxed{
P(Y=y\mid\theta_j).
}
\]

These describe how likely each signal is under each state of nature.

%%%%%%%%%%%%%%%%%%%%%%%%%%%%%%%%%%%%%%%%%%%%%%%%%%%%%%%%%%%%
\subsection{Bayes' Theorem}
\label{subsec:bayes-theorem-decision}
%%%%%%%%%%%%%%%%%%%%%%%%%%%%%%%%%%%%%%%%%%%%%%%%%%%%%%%%%%%%

After observing \(Y=y\), update beliefs using

\[
\boxed{
P(\theta_j\mid y)
=
\frac{
P(y\mid\theta_j)
P(\theta_j)
}{
P(y)
}.
}
\]

The denominator is

\[
\boxed{
P(y)
=
\sum_k
P(y\mid\theta_k)
P(\theta_k).
}
\]

%%%%%%%%%%%%%%%%%%%%%%%%%%%%%%%%%%%%%%%%%%%%%%%%%%%%%%%%%%%%
\subsection{Posterior Probabilities}
\label{subsec:posterior-probabilities}
%%%%%%%%%%%%%%%%%%%%%%%%%%%%%%%%%%%%%%%%%%%%%%%%%%%%%%%%%%%%

The probabilities

\[
\boxed{
P(\theta_j\mid y)
}
\]

are posterior probabilities.

They summarize the decision maker's updated beliefs after observing the
signal.

The optimal action may therefore depend on \(y\).

%%%%%%%%%%%%%%%%%%%%%%%%%%%%%%%%%%%%%%%%%%%%%%%%%%%%%%%%%%%%
\subsection{Worked Example: Bayesian Updating}
\label{subsec:worked-bayesian-updating}
%%%%%%%%%%%%%%%%%%%%%%%%%%%%%%%%%%%%%%%%%%%%%%%%%%%%%%%%%%%%

\begin{workedexamplebox}{Updating Demand Beliefs}

Suppose:

\[
P(H)=0.4,
\qquad
P(L)=0.6.
\]

A favorable market signal \(F\) satisfies

\[
P(F\mid H)=0.8,
\]

and

\[
P(F\mid L)=0.3.
\]

Then

\[
P(F)
=
0.8(0.4)
+
0.3(0.6)
=
0.5.
\]

Therefore,

\[
P(H\mid F)
=
\frac{
0.8(0.4)
}{
0.5
}
=
0.64.
\]

Similarly,

\[
P(L\mid F)
=
0.36.
\]

\begin{answerbox}
\[
\boxed{
P(H\mid F)=0.64,
\qquad
P(L\mid F)=0.36.
}
\]
\end{answerbox}

\end{workedexamplebox}

%%%%%%%%%%%%%%%%%%%%%%%%%%%%%%%%%%%%%%%%%%%%%%%%%%%%%%%%%%%%
\subsection{Expected Value of Sample Information}
\label{subsec:evsi}
%%%%%%%%%%%%%%%%%%%%%%%%%%%%%%%%%%%%%%%%%%%%%%%%%%%%%%%%%%%%

Let

\[
EV_{\mathrm{SI}}
\]

be the expected value obtained when sample information is observed before
the decision.

Then

\[
\boxed{
EVSI
=
EV_{\mathrm{SI}}
-
EV_{\mathrm{current}}.
}
\]

Because imperfect information cannot be more valuable than perfect
information,

\[
\boxed{
0
\leq
EVSI
\leq
EVPI.
}
\]

%%%%%%%%%%%%%%%%%%%%%%%%%%%%%%%%%%%%%%%%%%%%%%%%%%%%%%%%%%%%
\subsection{Net Value of Information}
\label{subsec:net-value-information}
%%%%%%%%%%%%%%%%%%%%%%%%%%%%%%%%%%%%%%%%%%%%%%%%%%%%%%%%%%%%

Suppose information costs

\[
C_I.
\]

Its net expected value is

\[
\boxed{
EVSI-C_I.
}
\]

Purchasing information is attractive under expected-value reasoning when

\[
\boxed{
EVSI>C_I.
}
\]

%%%%%%%%%%%%%%%%%%%%%%%%%%%%%%%%%%%%%%%%%%%%%%%%%%%%%%%%%%%%
\subsection{Information Has Decision Value Only If It Can Change Action}
\label{subsec:information-change-action}
%%%%%%%%%%%%%%%%%%%%%%%%%%%%%%%%%%%%%%%%%%%%%%%%%%%%%%%%%%%%

Information can improve a decision only if different observations can
change which action is optimal or improve how actions are chosen.

A highly accurate signal may have little value if the same action remains
optimal under every reasonable posterior belief.

%%%%%%%%%%%%%%%%%%%%%%%%%%%%%%%%%%%%%%%%%%%%%%%%%%%%%%%%%%%%
\subsection{Bayesian Decision Theory}
\label{subsec:bayesian-decision-theory}
%%%%%%%%%%%%%%%%%%%%%%%%%%%%%%%%%%%%%%%%%%%%%%%%%%%%%%%%%%%%

Bayesian decision theory combines:

\[
\boxed{
\text{prior probabilities},
}
\]

\[
\boxed{
\text{data},
}
\]

\[
\boxed{
\text{posterior probabilities},
}
\]

and

\[
\boxed{
\text{loss functions}.
}
\]

After observing data \(y\), choose action \(a\) minimizing posterior
expected loss.

%%%%%%%%%%%%%%%%%%%%%%%%%%%%%%%%%%%%%%%%%%%%%%%%%%%%%%%%%%%%
\subsection{Posterior Expected Loss}
\label{subsec:posterior-expected-loss}
%%%%%%%%%%%%%%%%%%%%%%%%%%%%%%%%%%%%%%%%%%%%%%%%%%%%%%%%%%%%

Given observation \(y\),

\[
\boxed{
\rho(a\mid y)
=
\mathbb{E}
[
L(a,\Theta)
\mid
Y=y
].
}
\]

For discrete states,

\[
\boxed{
\rho(a\mid y)
=
\sum_j
L(a,\theta_j)
P(\theta_j\mid y).
}
\]

The Bayes action is

\[
\boxed{
a^*(y)
\in
\operatorname*{arg\,min}_a
\rho(a\mid y).
}
\]

%%%%%%%%%%%%%%%%%%%%%%%%%%%%%%%%%%%%%%%%%%%%%%%%%%%%%%%%%%%%
\subsection{Bayes Risk}
\label{subsec:bayes-risk}
%%%%%%%%%%%%%%%%%%%%%%%%%%%%%%%%%%%%%%%%%%%%%%%%%%%%%%%%%%%%

The Bayes risk of a decision rule \(\delta\) is its expected loss under the
prior distribution and sampling model:

\[
\boxed{
r(\pi,\delta)
=
\mathbb{E}
[
L
(
\delta(Y),
\Theta
)
].
}
\]

A Bayes rule minimizes this expected loss.

%%%%%%%%%%%%%%%%%%%%%%%%%%%%%%%%%%%%%%%%%%%%%%%%%%%%%%%%%%%%
\subsection{Statistical Decision Problem}
\label{subsec:statistical-decision-problem}
%%%%%%%%%%%%%%%%%%%%%%%%%%%%%%%%%%%%%%%%%%%%%%%%%%%%%%%%%%%%

Statistical inference can be viewed as decision theory.

The unknown parameter is the state of nature:

\[
\boxed{
\theta.
}
\]

Observed data are

\[
\boxed{
X.
}
\]

The statistical procedure chooses an action

\[
\boxed{
a=\delta(X).
}
\]

A loss function determines the cost of estimation or classification error.

%%%%%%%%%%%%%%%%%%%%%%%%%%%%%%%%%%%%%%%%%%%%%%%%%%%%%%%%%%%%
\subsection{Squared-Error Loss}
\label{subsec:squared-error-loss}
%%%%%%%%%%%%%%%%%%%%%%%%%%%%%%%%%%%%%%%%%%%%%%%%%%%%%%%%%%%%

For estimating parameter \(\theta\), squared-error loss is

\[
\boxed{
L(a,\theta)
=
(a-\theta)^2.
}
\]

Under posterior squared-error loss, the Bayes estimator is the posterior
mean:

\[
\boxed{
a^*
=
\mathbb{E}
[
\Theta\mid X
].
}
\]

%%%%%%%%%%%%%%%%%%%%%%%%%%%%%%%%%%%%%%%%%%%%%%%%%%%%%%%%%%%%
\subsection{Absolute-Error Loss}
\label{subsec:absolute-error-loss}
%%%%%%%%%%%%%%%%%%%%%%%%%%%%%%%%%%%%%%%%%%%%%%%%%%%%%%%%%%%%

Absolute-error loss is

\[
\boxed{
L(a,\theta)
=
|a-\theta|.
}
\]

A posterior median minimizes posterior expected absolute error.

Thus the optimal estimator depends on the chosen loss function.

%%%%%%%%%%%%%%%%%%%%%%%%%%%%%%%%%%%%%%%%%%%%%%%%%%%%%%%%%%%%
\subsection{Zero-One Loss}
\label{subsec:zero-one-loss}
%%%%%%%%%%%%%%%%%%%%%%%%%%%%%%%%%%%%%%%%%%%%%%%%%%%%%%%%%%%%

For classification or discrete-state decisions, zero-one loss is

\[
\boxed{
L(a,\theta)
=
\begin{cases}
0,
&
a=\theta,
\\
1,
&
a\neq\theta.
\end{cases}
}
\]

Under zero-one loss, the Bayes action selects a state with maximum posterior
probability.

%%%%%%%%%%%%%%%%%%%%%%%%%%%%%%%%%%%%%%%%%%%%%%%%%%%%%%%%%%%%
\subsection{MAP Decision Rule}
\label{subsec:map-decision-rule}
%%%%%%%%%%%%%%%%%%%%%%%%%%%%%%%%%%%%%%%%%%%%%%%%%%%%%%%%%%%%

The maximum a posteriori rule chooses

\[
\boxed{
a^*
\in
\operatorname*{arg\,max}_{\theta_j}
P(\theta_j\mid y).
}
\]

This is optimal under zero-one loss.

With unequal misclassification costs, the MAP rule may no longer be
optimal.

%%%%%%%%%%%%%%%%%%%%%%%%%%%%%%%%%%%%%%%%%%%%%%%%%%%%%%%%%%%%
\subsection{Binary Classification as Decision Theory}
\label{subsec:binary-classification-decision}
%%%%%%%%%%%%%%%%%%%%%%%%%%%%%%%%%%%%%%%%%%%%%%%%%%%%%%%%%%%%

Suppose the true state is

\[
\Theta\in\{0,1\},
\]

and the decision is

\[
a\in\{0,1\}.
\]

Possible outcomes include:

\[
\boxed{
\text{true positive},
}
\]

\[
\boxed{
\text{true negative},
}
\]

\[
\boxed{
\text{false positive},
}
\]

and

\[
\boxed{
\text{false negative}.
}
\]

Different errors can have different costs.

%%%%%%%%%%%%%%%%%%%%%%%%%%%%%%%%%%%%%%%%%%%%%%%%%%%%%%%%%%%%
\subsection{Unequal Error Costs}
\label{subsec:unequal-error-costs}
%%%%%%%%%%%%%%%%%%%%%%%%%%%%%%%%%%%%%%%%%%%%%%%%%%%%%%%%%%%%

Suppose false-positive loss is

\[
C_{FP},
\]

and false-negative loss is

\[
C_{FN}.
\]

Then a classification decision should account for these losses rather than
using probability \(0.5\) automatically as the decision threshold.

%%%%%%%%%%%%%%%%%%%%%%%%%%%%%%%%%%%%%%%%%%%%%%%%%%%%%%%%%%%%
\subsection{Optimal Binary Decision Threshold}
\label{subsec:optimal-binary-threshold}
%%%%%%%%%%%%%%%%%%%%%%%%%%%%%%%%%%%%%%%%%%%%%%%%%%%%%%%%%%%%

Suppose correct classifications have zero loss.

Choose positive action \(a=1\) if

\[
C_{FP}
P(\Theta=0\mid y)
<
C_{FN}
P(\Theta=1\mid y).
\]

Equivalently,

\[
\boxed{
\frac{
P(\Theta=1\mid y)
}{
P(\Theta=0\mid y)
}
>
\frac{
C_{FP}
}{
C_{FN}
}.
}
\]

Thus loss asymmetry changes the optimal threshold.

%%%%%%%%%%%%%%%%%%%%%%%%%%%%%%%%%%%%%%%%%%%%%%%%%%%%%%%%%%%%
\subsection{Worked Example: Classification Decision}
\label{subsec:worked-classification-decision}
%%%%%%%%%%%%%%%%%%%%%%%%%%%%%%%%%%%%%%%%%%%%%%%%%%%%%%%%%%%%

\begin{workedexamplebox}{Unequal Misclassification Costs}

Suppose

\[
P(\Theta=1\mid y)=0.3,
\]

and

\[
P(\Theta=0\mid y)=0.7.
\]

Suppose false-positive cost is

\[
C_{FP}=1,
\]

while false-negative cost is

\[
C_{FN}=5.
\]

Expected loss from predicting positive is

\[
1(0.7)
=
0.7.
\]

Expected loss from predicting negative is

\[
5(0.3)
=
1.5.
\]

Therefore,

\begin{answerbox}
\[
\boxed{
\text{Predict positive}.
}
\]
\end{answerbox}

Even though the posterior probability of the positive state is less than
\(0.5\), the high false-negative cost changes the optimal action.

\end{workedexamplebox}

%%%%%%%%%%%%%%%%%%%%%%%%%%%%%%%%%%%%%%%%%%%%%%%%%%%%%%%%%%%%
\subsection{Sensitivity and Specificity}
\label{subsec:sensitivity-specificity-decision}
%%%%%%%%%%%%%%%%%%%%%%%%%%%%%%%%%%%%%%%%%%%%%%%%%%%%%%%%%%%%

For binary detection problems:

\[
\boxed{
\text{sensitivity}
=
P
(
\text{positive decision}
\mid
\text{positive state}
),
}
\]

and

\[
\boxed{
\text{specificity}
=
P
(
\text{negative decision}
\mid
\text{negative state}
).
}
\]

Changing the decision threshold typically changes both measures.

%%%%%%%%%%%%%%%%%%%%%%%%%%%%%%%%%%%%%%%%%%%%%%%%%%%%%%%%%%%%
\subsection{ROC Tradeoff}
\label{subsec:roc-decision-theory}
%%%%%%%%%%%%%%%%%%%%%%%%%%%%%%%%%%%%%%%%%%%%%%%%%%%%%%%%%%%%

A receiver operating characteristic, or ROC, curve studies the tradeoff
between:

\[
\boxed{
\text{true-positive rate}
}
\]

and

\[
\boxed{
\text{false-positive rate}.
}
\]

Decision theory determines where on this tradeoff curve an action should be
chosen after incorporating probabilities and losses.

%%%%%%%%%%%%%%%%%%%%%%%%%%%%%%%%%%%%%%%%%%%%%%%%%%%%%%%%%%%%
\subsection{Stochastic Dominance and Decisions}
\label{subsec:stochastic-dominance-decisions}
%%%%%%%%%%%%%%%%%%%%%%%%%%%%%%%%%%%%%%%%%%%%%%%%%%%%%%%%%%%%

If one risky alternative first-order stochastically dominates another,
then every decision maker with increasing utility prefers the dominant
alternative.

Thus stochastic dominance can eliminate some choices without requiring an
exact utility function.

%%%%%%%%%%%%%%%%%%%%%%%%%%%%%%%%%%%%%%%%%%%%%%%%%%%%%%%%%%%%
\subsection{Decision Making With Multiple Attributes}
\label{subsec:multiattribute-decision-making}
%%%%%%%%%%%%%%%%%%%%%%%%%%%%%%%%%%%%%%%%%%%%%%%%%%%%%%%%%%%%

Real decisions may involve outcomes described by several attributes:

\[
\boxed{
x
=
(x_1,\ldots,x_k).
}
\]

Examples include:

\[
\boxed{
\text{cost},
}
\]

\[
\boxed{
\text{time},
}
\]

\[
\boxed{
\text{safety},
}
\]

and

\[
\boxed{
\text{quality}.
}
\]

A utility function can combine these attributes.

%%%%%%%%%%%%%%%%%%%%%%%%%%%%%%%%%%%%%%%%%%%%%%%%%%%%%%%%%%%%
\subsection{Additive Multiattribute Utility}
\label{subsec:additive-multiattribute-utility}
%%%%%%%%%%%%%%%%%%%%%%%%%%%%%%%%%%%%%%%%%%%%%%%%%%%%%%%%%%%%

Under suitable preference-independence conditions, one may use

\[
\boxed{
U(x)
=
\sum_{i=1}^{k}
w_i
u_i(x_i),
}
\]

where

\[
w_i\geq0,
\]

and often

\[
\boxed{
\sum_iw_i=1.
}
\]

This resembles weighted-sum multiobjective optimization.

%%%%%%%%%%%%%%%%%%%%%%%%%%%%%%%%%%%%%%%%%%%%%%%%%%%%%%%%%%%%
\subsection{Utility Versus Objective Weights}
\label{subsec:utility-versus-objective-weights}
%%%%%%%%%%%%%%%%%%%%%%%%%%%%%%%%%%%%%%%%%%%%%%%%%%%%%%%%%%%%

Objective weights in multiobjective optimization and utility coefficients
both encode preferences, but they arise from different modeling
frameworks.

Utility theory attempts to model preferences over consequences and risk.

Weighted optimization combines mathematical objective functions into a
single scalar criterion.

Care is needed when interpreting either type of weight.

%%%%%%%%%%%%%%%%%%%%%%%%%%%%%%%%%%%%%%%%%%%%%%%%%%%%%%%%%%%%
\subsection{Preference Elicitation}
\label{subsec:preference-elicitation}
%%%%%%%%%%%%%%%%%%%%%%%%%%%%%%%%%%%%%%%%%%%%%%%%%%%%%%%%%%%%

Utility or decision-model parameters may be elicited using:

\[
\boxed{
\text{tradeoff questions},
}
\]

\[
\boxed{
\text{certainty-equivalent questions},
}
\]

\[
\boxed{
\text{indifference judgments},
}
\]

or

\[
\boxed{
\text{observed choices}.
}
\]

Poorly elicited preferences can make a mathematically correct model
decision practically misleading.

%%%%%%%%%%%%%%%%%%%%%%%%%%%%%%%%%%%%%%%%%%%%%%%%%%%%%%%%%%%%
\subsection{Sensitivity Analysis}
\label{subsec:decision-sensitivity-analysis}
%%%%%%%%%%%%%%%%%%%%%%%%%%%%%%%%%%%%%%%%%%%%%%%%%%%%%%%%%%%%

Decision recommendations may depend strongly on:

\[
\boxed{
\text{probabilities},
}
\]

\[
\boxed{
\text{payoffs},
}
\]

\[
\boxed{
\text{losses},
}
\]

or

\[
\boxed{
\text{utility parameters}.
}
\]

Sensitivity analysis asks whether the chosen action changes when these
inputs vary.

%%%%%%%%%%%%%%%%%%%%%%%%%%%%%%%%%%%%%%%%%%%%%%%%%%%%%%%%%%%%
\subsection{Threshold Probability Analysis}
\label{subsec:threshold-probability-analysis}
%%%%%%%%%%%%%%%%%%%%%%%%%%%%%%%%%%%%%%%%%%%%%%%%%%%%%%%%%%%%

Suppose state \(\theta_1\) has probability \(p\), and state \(\theta_2\)
has probability \(1-p\).

For two actions \(A\) and \(B\),

\[
EMV(A)
=
pu_{A1}
+
(1-p)u_{A2},
\]

and

\[
EMV(B)
=
pu_{B1}
+
(1-p)u_{B2}.
\]

Set them equal to find the probability threshold at which the preferred
action changes.

%%%%%%%%%%%%%%%%%%%%%%%%%%%%%%%%%%%%%%%%%%%%%%%%%%%%%%%%%%%%
\subsection{Worked Example: Probability Threshold}
\label{subsec:worked-probability-threshold}
%%%%%%%%%%%%%%%%%%%%%%%%%%%%%%%%%%%%%%%%%%%%%%%%%%%%%%%%%%%%

\begin{workedexamplebox}{When Does the Preferred Plan Change?}

Suppose

\[
A:
(100,20),
\]

and

\[
B:
(70,60).
\]

Let \(p\) be the probability of the first state.

Then

\[
EMV(A)
=
100p
+
20(1-p)
=
20+80p.
\]

Also,

\[
EMV(B)
=
70p
+
60(1-p)
=
60+10p.
\]

Set equal:

\[
20+80p
=
60+10p.
\]

Thus,

\[
70p
=
40,
\]

so

\[
p
=
\frac47.
\]

Therefore,

\begin{answerbox}
\[
\boxed{
p^*
=
\frac47
\approx0.5714.
}
\]
\end{answerbox}

Above this threshold, \(A\) has larger expected payoff.

Below it, \(B\) does.

\end{workedexamplebox}

%%%%%%%%%%%%%%%%%%%%%%%%%%%%%%%%%%%%%%%%%%%%%%%%%%%%%%%%%%%%
\subsection{Decision Theory and Dynamic Programming}
\label{subsec:decision-theory-dynamic-programming}
%%%%%%%%%%%%%%%%%%%%%%%%%%%%%%%%%%%%%%%%%%%%%%%%%%%%%%%%%%%%

Sequential decision problems can be solved recursively.

At each stage, one:

\[
\boxed{
\text{chooses an action},
}
\]

observes uncertainty, and then

\[
\boxed{
\text{chooses again}.
}
\]

This leads naturally to Bellman equations.

Decision trees are finite representations of the same principle.

%%%%%%%%%%%%%%%%%%%%%%%%%%%%%%%%%%%%%%%%%%%%%%%%%%%%%%%%%%%%
\subsection{Decision Theory and Stochastic Optimization}
\label{subsec:decision-theory-stochastic-optimization}
%%%%%%%%%%%%%%%%%%%%%%%%%%%%%%%%%%%%%%%%%%%%%%%%%%%%%%%%%%%%

Stochastic optimization focuses on finding optimal decisions in models
containing random variables.

Decision theory adds explicit attention to:

\[
\boxed{
\text{preferences},
}
\]

\[
\boxed{
\text{information},
}
\]

and

\[
\boxed{
\text{loss or utility}.
}
\]

The fields overlap substantially.

%%%%%%%%%%%%%%%%%%%%%%%%%%%%%%%%%%%%%%%%%%%%%%%%%%%%%%%%%%%%
\subsection{Decision Theory and Game Theory}
\label{subsec:decision-theory-game-theory}
%%%%%%%%%%%%%%%%%%%%%%%%%%%%%%%%%%%%%%%%%%%%%%%%%%%%%%%%%%%%

In classical decision theory, uncertainty often comes from a state of
nature that does not strategically respond to the decision maker.

In game theory, outcomes depend on choices made by other strategic agents.

Thus:

\[
\boxed{
\text{decision theory}
\rightarrow
\text{uncertain environment},
}
\]

while

\[
\boxed{
\text{game theory}
\rightarrow
\text{strategic interaction}.
}
\]

%%%%%%%%%%%%%%%%%%%%%%%%%%%%%%%%%%%%%%%%%%%%%%%%%%%%%%%%%%%%
\subsection{Ambiguity}
\label{subsec:decision-ambiguity}
%%%%%%%%%%%%%%%%%%%%%%%%%%%%%%%%%%%%%%%%%%%%%%%%%%%%%%%%%%%%

Risk refers to uncertainty described by known probabilities.

Ambiguity refers to uncertainty about the probability model itself.

A decision maker may know possible states but not know reliable values of

\[
\boxed{
P(\theta_j).
}
\]

This motivates robust and ambiguity-sensitive approaches.

%%%%%%%%%%%%%%%%%%%%%%%%%%%%%%%%%%%%%%%%%%%%%%%%%%%%%%%%%%%%
\subsection{Robust Decision Making}
\label{subsec:robust-decision-making}
%%%%%%%%%%%%%%%%%%%%%%%%%%%%%%%%%%%%%%%%%%%%%%%%%%%%%%%%%%%%

Suppose the probability model belongs to a set

\[
\mathcal{P}.
\]

A robust expected-loss rule might solve

\[
\boxed{
\min_a
\sup_{P\in\mathcal{P}}
\mathbb{E}_{P}
[
L(a,\Theta)
].
}
\]

This protects against unfavorable distributions within the ambiguity set.

%%%%%%%%%%%%%%%%%%%%%%%%%%%%%%%%%%%%%%%%%%%%%%%%%%%%%%%%%%%%
\subsection{Minimax Decision Rule}
\label{subsec:minimax-decision-rule}
%%%%%%%%%%%%%%%%%%%%%%%%%%%%%%%%%%%%%%%%%%%%%%%%%%%%%%%%%%%%

Without a prior probability distribution, a minimax rule may choose

\[
\boxed{
a^*
\in
\operatorname*{arg\,min}_a
\sup_{\theta}
L(a,\theta).
}
\]

This minimizes worst-case loss.

It is the loss-version analogue of conservative decision criteria.

%%%%%%%%%%%%%%%%%%%%%%%%%%%%%%%%%%%%%%%%%%%%%%%%%%%%%%%%%%%%
\subsection{Minimax Regret and Robustness}
\label{subsec:minimax-regret-robustness}
%%%%%%%%%%%%%%%%%%%%%%%%%%%%%%%%%%%%%%%%%%%%%%%%%%%%%%%%%%%%

Worst-case loss and worst-case regret are not the same.

Minimax loss asks:

\[
\boxed{
\text{How bad can my loss become?}
}
\]

Minimax regret asks:

\[
\boxed{
\text{How much worse can I be than the action I would have chosen with
hindsight?}
}
\]

The selected actions may differ.

%%%%%%%%%%%%%%%%%%%%%%%%%%%%%%%%%%%%%%%%%%%%%%%%%%%%%%%%%%%%
\subsection{Behavioral Decision Theory Preview}
\label{subsec:behavioral-decision-theory}
%%%%%%%%%%%%%%%%%%%%%%%%%%%%%%%%%%%%%%%%%%%%%%%%%%%%%%%%%%%%

Observed human decisions often deviate from classical expected-utility
predictions.

Examples include:

\[
\boxed{
\text{loss aversion},
}
\]

\[
\boxed{
\text{probability weighting},
}
\]

\[
\boxed{
\text{framing effects},
}
\]

and

\[
\boxed{
\text{reference dependence}.
}
\]

Behavioral decision theory develops models of these patterns.

%%%%%%%%%%%%%%%%%%%%%%%%%%%%%%%%%%%%%%%%%%%%%%%%%%%%%%%%%%%%
\subsection{Prospect Theory Preview}
\label{subsec:prospect-theory-preview}
%%%%%%%%%%%%%%%%%%%%%%%%%%%%%%%%%%%%%%%%%%%%%%%%%%%%%%%%%%%%

Prospect theory models outcomes relative to a reference point.

Typical features include:

\[
\boxed{
\text{greater sensitivity to losses than comparable gains},
}
\]

and

\[
\boxed{
\text{nonlinear probability weighting}.
}
\]

This differs from standard expected utility over final wealth.

%%%%%%%%%%%%%%%%%%%%%%%%%%%%%%%%%%%%%%%%%%%%%%%%%%%%%%%%%%%%
\subsection{Decision Biases}
\label{subsec:decision-biases}
%%%%%%%%%%%%%%%%%%%%%%%%%%%%%%%%%%%%%%%%%%%%%%%%%%%%%%%%%%%%

Decision processes may be influenced by:

\[
\boxed{
\text{anchoring},
}
\]

\[
\boxed{
\text{availability effects},
}
\]

\[
\boxed{
\text{confirmation bias},
}
\]

or

\[
\boxed{
\text{overconfidence}.
}
\]

Mathematical decision analysis can clarify assumptions but does not
automatically eliminate human judgment errors.

%%%%%%%%%%%%%%%%%%%%%%%%%%%%%%%%%%%%%%%%%%%%%%%%%%%%%%%%%%%%
\subsection{Structured Decision Analysis}
\label{subsec:structured-decision-analysis}
%%%%%%%%%%%%%%%%%%%%%%%%%%%%%%%%%%%%%%%%%%%%%%%%%%%%%%%%%%%%

A structured decision process may proceed as follows:

\begin{enumerate}
    \item define the decision;
    \item identify alternatives;
    \item identify uncertain states;
    \item define consequences;
    \item estimate probabilities if possible;
    \item define losses or utility;
    \item identify available information;
    \item compute decision criteria;
    \item perform sensitivity analysis;
    \item evaluate robustness;
    \item choose and document the recommended action.
\end{enumerate}

%%%%%%%%%%%%%%%%%%%%%%%%%%%%%%%%%%%%%%%%%%%%%%%%%%%%%%%%%%%%
\subsection{Decision-Theory Workflow}
\label{subsec:decision-theory-workflow}
%%%%%%%%%%%%%%%%%%%%%%%%%%%%%%%%%%%%%%%%%%%%%%%%%%%%%%%%%%%%

A practical workflow is:

\begin{enumerate}
    \item state the decision clearly;
    \item list feasible actions;
    \item identify uncertain states of nature;
    \item construct a payoff or loss table;
    \item determine whether probabilities are available;
    \item if probabilities are available, compute expected values or
          expected utilities;
    \item if probabilities are unavailable, identify an appropriate
          uncertainty criterion;
    \item evaluate risk preferences;
    \item build a decision tree for sequential problems;
    \item update probabilities when information is observed;
    \item calculate the value of information;
    \item examine sensitivity to probabilities and consequences;
    \item identify dominated decisions;
    \item test robustness to model assumptions;
    \item document the final tradeoffs and assumptions.
\end{enumerate}

%%%%%%%%%%%%%%%%%%%%%%%%%%%%%%%%%%%%%%%%%%%%%%%%%%%%%%%%%%%%
\subsection{Choosing a Decision Criterion}
\label{subsec:choosing-decision-criterion}
%%%%%%%%%%%%%%%%%%%%%%%%%%%%%%%%%%%%%%%%%%%%%%%%%%%%%%%%%%%%

For known probabilities and approximately risk-neutral monetary outcomes,
consider:

\[
\boxed{
\text{expected monetary value}.
}
\]

For explicit risk preferences, consider:

\[
\boxed{
\text{expected utility}.
}
\]

For conservative decisions without probabilities, consider:

\[
\boxed{
\text{maximin or minimax loss}.
}
\]

For concern about hindsight opportunity loss, consider:

\[
\boxed{
\text{minimax regret}.
}
\]

For sequential decisions with information, consider:

\[
\boxed{
\text{decision trees or dynamic programming}.
}
\]

%%%%%%%%%%%%%%%%%%%%%%%%%%%%%%%%%%%%%%%%%%%%%%%%%%%%%%%%%%%%
\subsection{Common Errors}
\label{subsec:decision-theory-common-errors}
%%%%%%%%%%%%%%%%%%%%%%%%%%%%%%%%%%%%%%%%%%%%%%%%%%%%%%%%%%%%

\subsubsection{Confusing Decisions With States of Nature}

A decision is controlled by the decision maker.

A state of nature is not.

The two should occupy different dimensions of a decision table.

%%%%%%%%%%%%%%%%%%%%%%%%%%%%%%%%%%%%%%%%%%%%%%%%%%%%%%%%%%%%
\subsubsection{Mixing Payoffs and Losses}

For payoffs, larger is generally better.

For losses, smaller is generally better.

Using a maximization criterion on a loss table reverses the intended
decision.

%%%%%%%%%%%%%%%%%%%%%%%%%%%%%%%%%%%%%%%%%%%%%%%%%%%%%%%%%%%%
\subsubsection{Using Expected Value Without Probabilities}

Expected monetary value requires probabilities.

If probabilities are unavailable, assigning them implicitly without
justification changes the model.

%%%%%%%%%%%%%%%%%%%%%%%%%%%%%%%%%%%%%%%%%%%%%%%%%%%%%%%%%%%%
\subsubsection{Treating Equal Probabilities as Objective Ignorance}

The Laplace principle assigns equal probabilities as a decision rule.

It does not prove that equally likely states are objectively correct.

%%%%%%%%%%%%%%%%%%%%%%%%%%%%%%%%%%%%%%%%%%%%%%%%%%%%%%%%%%%%
\subsubsection{Confusing Maximin and Minimax Regret}

Maximin examines worst actual payoff.

Minimax regret examines worst opportunity loss.

They are different criteria.

%%%%%%%%%%%%%%%%%%%%%%%%%%%%%%%%%%%%%%%%%%%%%%%%%%%%%%%%%%%%
\subsubsection{Assuming Maximum Expected Money Is Always Rational}

Expected monetary value corresponds to risk neutrality.

A risk-averse or risk-seeking decision maker may prefer another action.

%%%%%%%%%%%%%%%%%%%%%%%%%%%%%%%%%%%%%%%%%%%%%%%%%%%%%%%%%%%%
\subsubsection{Confusing Expected Utility and Utility of Expected Value}

In general,

\[
\boxed{
\mathbb{E}[U(X)]
\neq
U(\mathbb{E}[X]).
}
\]

This distinction is central to expected-utility theory.

%%%%%%%%%%%%%%%%%%%%%%%%%%%%%%%%%%%%%%%%%%%%%%%%%%%%%%%%%%%%
\subsubsection{Using Utility Values as Dollar Amounts}

Utility is a preference representation.

It does not generally have the same units or interpretation as money.

%%%%%%%%%%%%%%%%%%%%%%%%%%%%%%%%%%%%%%%%%%%%%%%%%%%%%%%%%%%%
\subsubsection{Using the Wrong Sign for Risk Premium}

For a risk-averse decision maker,

\[
\boxed{
\pi
=
\mathbb{E}[W]-CE
\geq0.
}
\]

Reversing the terms changes its interpretation.

%%%%%%%%%%%%%%%%%%%%%%%%%%%%%%%%%%%%%%%%%%%%%%%%%%%%%%%%%%%%
\subsubsection{Ignoring Information Costs}

Information may have positive value but still not be worth purchasing if
its cost exceeds EVSI.

%%%%%%%%%%%%%%%%%%%%%%%%%%%%%%%%%%%%%%%%%%%%%%%%%%%%%%%%%%%%
\subsubsection{Assuming Imperfect Information Can Be More Valuable Than Perfect Information}

Under consistent modeling,

\[
\boxed{
EVSI
\leq
EVPI.
}
\]

Perfect information provides the strongest possible state revelation.

%%%%%%%%%%%%%%%%%%%%%%%%%%%%%%%%%%%%%%%%%%%%%%%%%%%%%%%%%%%%
\subsubsection{Using Prior Probabilities After Observing Data}

Once information is observed, decisions should use posterior probabilities
when following Bayesian decision analysis.

%%%%%%%%%%%%%%%%%%%%%%%%%%%%%%%%%%%%%%%%%%%%%%%%%%%%%%%%%%%%
\subsubsection{Ignoring Unequal Error Costs in Classification}

A probability threshold of \(0.5\) is not automatically optimal.

Losses and prior probabilities determine the optimal decision rule.

%%%%%%%%%%%%%%%%%%%%%%%%%%%%%%%%%%%%%%%%%%%%%%%%%%%%%%%%%%%%
\subsubsection{Assuming the Most Likely State Determines the Best Action}

The most probable state need not determine the optimal action.

Consequences and losses across all states matter.

%%%%%%%%%%%%%%%%%%%%%%%%%%%%%%%%%%%%%%%%%%%%%%%%%%%%%%%%%%%%
\subsubsection{Ignoring Sensitivity}

A recommendation may reverse under small probability or payoff changes.

Threshold analysis can reveal fragile decisions.

%%%%%%%%%%%%%%%%%%%%%%%%%%%%%%%%%%%%%%%%%%%%%%%%%%%%%%%%%%%%
\subsubsection{Confusing Nature With a Strategic Opponent}

A state of nature does not usually react strategically.

If another agent deliberately changes behavior in response to the
decision, game theory is more appropriate.

%%%%%%%%%%%%%%%%%%%%%%%%%%%%%%%%%%%%%%%%%%%%%%%%%%%%%%%%%%%%
\subsection{Applications}
\label{subsec:decision-theory-applications}
%%%%%%%%%%%%%%%%%%%%%%%%%%%%%%%%%%%%%%%%%%%%%%%%%%%%%%%%%%%%

\begin{applicationbox}

\textbf{Business Strategy.}
Organizations choose among investments, facilities, products, and
contracts before market conditions are known.

\medskip

\textbf{Healthcare.}
Treatment and diagnostic decisions balance probabilities, benefits,
risks, and asymmetric error costs.

\medskip

\textbf{Engineering.}
Design decisions must account for uncertain loads, failures, reliability,
and safety consequences.

\medskip

\textbf{Finance.}
Investment decisions combine uncertain returns with utility and risk
preferences.

\medskip

\textbf{Supply Chains.}
Capacity, sourcing, inventory, and contracting decisions depend on
uncertain demand and disruptions.

\medskip

\textbf{Public Policy.}
Policy decisions may involve uncertain outcomes, competing consequences,
and the value of additional information.

\medskip

\textbf{Environmental Management.}
Decisions may be made under uncertain climate, ecological, and economic
conditions.

\medskip

\textbf{Statistical Inference.}
Estimation and classification can be formulated as decisions minimizing
expected loss.

\end{applicationbox}

%%%%%%%%%%%%%%%%%%%%%%%%%%%%%%%%%%%%%%%%%%%%%%%%%%%%%%%%%%%%
\subsection{Exercises}
\label{subsec:decision-theory-exercises}
%%%%%%%%%%%%%%%%%%%%%%%%%%%%%%%%%%%%%%%%%%%%%%%%%%%%%%%%%%%%

\begin{exercise}[1]
Define a decision alternative.
\end{exercise}

\begin{exercise}[1]
Define a state of nature.
\end{exercise}

\begin{exercise}[1]
Define a payoff table.
\end{exercise}

\begin{exercise}[1]
Define regret.
\end{exercise}

\begin{exercise}[1]
Define expected monetary value.
\end{exercise}

\begin{exercise}[1]
Define utility.
\end{exercise}

\begin{exercise}[1]
Define certainty equivalent.
\end{exercise}

\begin{exercise}[1]
Define risk premium.
\end{exercise}

\begin{exercise}[1]
Define prior probability.
\end{exercise}

\begin{exercise}[1]
Define posterior probability.
\end{exercise}

\begin{exercise}[2]
Compute the EMV of an action with payoffs

\[
50,\ 80,\ 20
\]

and probabilities

\[
0.2,\ 0.5,\ 0.3.
\]
\end{exercise}

\begin{exercise}[2]
Use maximin to choose among a given set of payoff rows.
\end{exercise}

\begin{exercise}[2]
Use maximax on the same payoff table.
\end{exercise}

\begin{exercise}[2]
Apply the Hurwicz criterion with

\[
\alpha=0.7.
\]
\end{exercise}

\begin{exercise}[2]
Apply the Laplace principle to a three-state payoff table.
\end{exercise}

\begin{exercise}[2]
Construct a regret table from a given payoff table.
\end{exercise}

\begin{exercise}[2]
Use minimax regret to choose an action.
\end{exercise}

\begin{exercise}[3]
Compare the decisions obtained from maximin, maximax, Hurwicz, Laplace,
and minimax regret.
\end{exercise}

\begin{exercise}[3]
Explain why different uncertainty criteria may select different actions.
\end{exercise}

\begin{exercise}[3]
For a lottery with two monetary outcomes, compute expected wealth.
\end{exercise}

\begin{exercise}[3]
Using

\[
U(w)=\sqrt{w},
\]

compute expected utility.
\end{exercise}

\begin{exercise}[3]
Find the certainty equivalent of the previous lottery.
\end{exercise}

\begin{exercise}[3]
Compute the risk premium.
\end{exercise}

\begin{exercise}[3]
Use Jensen's inequality to explain risk aversion for concave utility.
\end{exercise}

\begin{exercise}[3]
Compute absolute risk aversion for

\[
U(w)=\ln w.
\]
\end{exercise}

\begin{exercise}[3]
Compute relative risk aversion for the same utility.
\end{exercise}

\begin{exercise}[3]
Determine whether one of two discrete lotteries first-order stochastically
dominates the other.
\end{exercise}

\begin{exercise}[3]
Construct a two-stage decision tree.
\end{exercise}

\begin{exercise}[3]
Evaluate the previous tree using rollback analysis.
\end{exercise}

\begin{exercise}[3]
Compute EVPI from a payoff table and state probabilities.
\end{exercise}

\begin{exercise}[3]
Explain why EVPI cannot be negative in a standard payoff-maximization
problem.
\end{exercise}

\begin{exercise}[3]
Given a diagnostic signal model, compute posterior probabilities using
Bayes' theorem.
\end{exercise}

\begin{exercise}[3]
Calculate the expected payoff after observing sample information.
\end{exercise}

\begin{exercise}[3]
Compute EVSI.
\end{exercise}

\begin{exercise}[3]
Verify that

\[
EVSI\leq EVPI.
\]
\end{exercise}

\begin{exercise}[3]
Determine whether information should be purchased given its cost.
\end{exercise}

\begin{exercise}[3]
Construct posterior expected-loss calculations for two actions.
\end{exercise}

\begin{exercise}[3]
Find the Bayes action.
\end{exercise}

\begin{exercise}[3]
Show that the posterior mean minimizes squared-error loss.
\end{exercise}

\begin{exercise}[3]
Explain why a posterior median minimizes absolute-error loss.
\end{exercise}

\begin{exercise}[3]
Show that the MAP estimate minimizes posterior zero-one loss.
\end{exercise}

\begin{exercise}[3]
Construct a binary classification loss matrix.
\end{exercise}

\begin{exercise}[3]
Determine the optimal binary action when false-positive and false-negative
costs differ.
\end{exercise}

\begin{exercise}[3]
Explain why a \(0.5\) classification threshold may be inappropriate.
\end{exercise}

\begin{exercise}[3]
Compute a probability threshold at which two decisions have equal EMV.
\end{exercise}

\begin{exercise}[3]
Perform sensitivity analysis around the threshold.
\end{exercise}

\begin{exercise}[3]
Construct a two-attribute additive utility model.
\end{exercise}

\begin{exercise}[3]
Explain the relationship between multiattribute utility and
multiobjective optimization.
\end{exercise}

\begin{exercise}[3]
Compare minimax loss and minimax regret on the same problem.
\end{exercise}

\begin{exercise}[3]
Explain the difference between risk and ambiguity.
\end{exercise}

\begin{exercise}[3]
Formulate a robust minimax expected-loss decision rule.
\end{exercise}

\begin{exercise}[4]
Study the von Neumann--Morgenstern expected-utility theorem.
\end{exercise}

\begin{exercise}[4]
Study affine uniqueness of expected-utility representations.
\end{exercise}

\begin{exercise}[4]
Derive local risk-premium approximations using Arrow--Pratt risk aversion.
\end{exercise}

\begin{exercise}[4]
Study constant absolute risk-aversion utility.
\end{exercise}

\begin{exercise}[4]
Study constant relative risk-aversion utility.
\end{exercise}

\begin{exercise}[4]
Prove basic first-order stochastic-dominance results.
\end{exercise}

\begin{exercise}[4]
Study second-order stochastic dominance and concave utility.
\end{exercise}

\begin{exercise}[4]
Study statistical decision theory with continuous parameter spaces.
\end{exercise}

\begin{exercise}[4]
Derive Bayes estimators under alternative loss functions.
\end{exercise}

\begin{exercise}[4]
Study admissible and inadmissible decision rules.
\end{exercise}

\begin{exercise}[4]
Study minimax statistical decision rules.
\end{exercise}

\begin{exercise}[4]
Investigate complete classes of decision rules.
\end{exercise}

\begin{exercise}[4]
Study likelihood-ratio decision rules for binary testing.
\end{exercise}

\begin{exercise}[4]
Derive Bayes classification thresholds from likelihood ratios.
\end{exercise}

\begin{exercise}[4]
Study ROC curves through decision-theoretic loss functions.
\end{exercise}

\begin{exercise}[4]
Study multiattribute utility theory.
\end{exercise}

\begin{exercise}[4]
Investigate value-of-information analysis for sequential decisions.
\end{exercise}

\begin{exercise}[4]
Study robust Bayes decision rules.
\end{exercise}

\begin{exercise}[4]
Study ambiguity-sensitive preferences.
\end{exercise}

\begin{exercise}[4]
Study prospect theory and compare it with expected utility.
\end{exercise}

%%%%%%%%%%%%%%%%%%%%%%%%%%%%%%%%%%%%%%%%%%%%%%%%%%%%%%%%%%%%
\subsection{Conceptual Summary}
\label{subsec:decision-theory-summary}
%%%%%%%%%%%%%%%%%%%%%%%%%%%%%%%%%%%%%%%%%%%%%%%%%%%%%%%%%%%%

Decision theory studies choices whose consequences depend on uncertain
states.

A basic decision model contains

\[
\boxed{
\text{actions}
+
\text{states}
+
\text{payoffs or losses}.
}
\]

Under risk, expected monetary value is

\[
\boxed{
EMV(a_i)
=
\sum_j
p_j
u_{ij}.
}
\]

Under uncertainty without probabilities, common criteria include

\[
\boxed{
\text{maximin},
}
\]

\[
\boxed{
\text{maximax},
}
\]

\[
\boxed{
\text{Hurwicz},
}
\]

and

\[
\boxed{
\text{minimax regret}.
}
\]

Expected utility is

\[
\boxed{
EU
=
\mathbb{E}[U(X)].
}
\]

The certainty equivalent satisfies

\[
\boxed{
U(CE)
=
\mathbb{E}[U(X)].
}
\]

The risk premium is

\[
\boxed{
\pi
=
\mathbb{E}[X]-CE.
}
\]

Bayesian decision theory updates

\[
\boxed{
\text{prior}
\rightarrow
\text{data}
\rightarrow
\text{posterior}
}
\]

and then chooses an action minimizing posterior expected loss.

The expected value of perfect information is

\[
\boxed{
EVPI
=
EV_{\mathrm{PI}}
-
EV_{\mathrm{current}}.
}
\]

For imperfect information,

\[
\boxed{
0
\leq
EVSI
\leq
EVPI.
}
\]

Decision theory therefore combines

\[
\boxed{
\text{probability}
+
\text{utility}
+
\text{statistics}
+
\text{optimization}
+
\text{information}.
}
\]

%%%%%%%%%%%%%%%%%%%%%%%%%%%%%%%%%%%%%%%%%%%%%%%%%%%%%%%%%%%%
\subsection{Section Connections}
\label{subsec:decision-theory-connections}
%%%%%%%%%%%%%%%%%%%%%%%%%%%%%%%%%%%%%%%%%%%%%%%%%%%%%%%%%%%%

\chapterconnection{
Decision Theory provides a mathematical framework for choosing among
alternatives when consequences depend on uncertainty, information, and
preferences. Expected-value analysis connects with
Section~\ref{sec:stochastic-optimization}, while multiattribute utility
connects with Section~\ref{sec:multiobjective-optimization}. Bayesian
decision rules connect directly with the probability and statistics
chapters, and sequential decision trees connect with
Section~\ref{sec:dynamic-programming}. The next section, Game Theory,
changes the source of uncertainty: rather than treating uncertain states
as nonstrategic nature, it studies situations in which outcomes depend on
the deliberate choices of multiple interacting decision makers.
}

%%%%%%%%%%%%%%%%%%%%%%%%%%%%%%%%%%%%%%%%%%%%%%%%%%%%%%%%%%%%
% SECTION SOURCE TRACKING
%%%%%%%%%%%%%%%%%%%%%%%%%%%%%%%%%%%%%%%%%%%%%%%%%%%%%%%%%%%%

%------------------------------------------------------------
% 10.14 Decision Theory
%------------------------------------------------------------
%
% Source basis:
%   - Standard classical and Bayesian decision theory.
%   - Standard expected-utility theory.
%   - Standard statistical decision analysis.
%
% Used for:
%   - actions and states of nature;
%   - payoff and loss tables;
%   - certainty, risk, and uncertainty;
%   - expected monetary value;
%   - expected loss;
%   - maximax and maximin;
%   - Hurwicz criterion;
%   - Laplace principle;
%   - regret and minimax regret;
%   - utility functions;
%   - expected utility;
%   - risk aversion, neutrality, and seeking;
%   - Jensen's inequality;
%   - certainty equivalents;
%   - risk premiums;
%   - Arrow--Pratt risk aversion;
%   - common utility functions;
%   - expected-utility axioms;
%   - stochastic dominance;
%   - decision trees;
%   - rollback analysis;
%   - perfect information;
%   - EVPI;
%   - sample information;
%   - Bayes' theorem;
%   - posterior probabilities;
%   - EVSI;
%   - Bayesian decision theory;
%   - posterior expected loss;
%   - Bayes risk;
%   - statistical decision theory;
%   - squared, absolute, and zero-one loss;
%   - MAP decisions;
%   - binary classification;
%   - asymmetric error costs;
%   - sensitivity and specificity;
%   - ROC tradeoffs;
%   - multiattribute utility;
%   - preference elicitation;
%   - sensitivity analysis;
%   - probability thresholds;
%   - robust and minimax decisions;
%   - ambiguity;
%   - behavioral decision-theory preview;
%   - applications and exercises.
%
% Notes:
%   - Game Theory is developed next in Section 10.15.
%   - Optimal Control follows in Section 10.16.
%   - Bayesian foundations are developed more deeply in Chapter 8.
%   - Stochastic Optimization appears in Section 10.8.
%
%%%%%%%%%%%%%%%%%%%%%%%%%%%%%%%%%%%%%%%%%%%%%%%%%%%%%%%%%%%%